
The philosophical analysis thus far might suggest designers are moral failures, choosing profit over planet. This interpretation is too simple. Designers operate within structural constraints that severely limit individual agency.

\subsection{Professional Pressure}
Consider the possible professional reality: Designers often encounter significant challenges when integrating environmental sustainability into their work, primarily due to constraints imposed by organisational priorities and policies \cite{CALABRETTA2025145340}. Designers are evaluated on engagement metrics, not sustainability. Proposing to disable autoplay or default to standard definition video quality might mean answering the question: "What impact will this have on engagement?" The honest answer "engagement will likely drop 10-15\%” could end the conversation. The proposal is rejected because the company's survival depends on growth metrics. Sustainability is a complex issue that involves multiple stakeholders, many not sharing the same level of enthusiasm for sustainability \cite{CALABRETTA2025145340}. A designer who consistently prioritizes environmental impact over business metrics might face performance reviews questioning their alignment with company goals or might ultimately be replacement by someone less resistant. The gap between personal values and anticipated professional behavior reveals the structural nature of the problem. Designers retain agency. They make choices and bear some responsibility. But agency operates within severe constraints determined by business models that require infinite growth.

\subsection{Interactive Exercise}
To examine how ethical frameworks apply to real design decisions, I conducted an interactive exercise with 5 design students during this papers presentation. Participants were presented with a scenario: as a UX designer for a streaming platform, choose between two default configurations. Option A (Industry Standard): 4K default with autoplay enabled, yielding +30\% engagement but 280g CO₂ per user per week. Option B (Planet-First): HD default with autoplay disabled, accepting -10\% engagement but generating only 45g CO₂ per user per week. Voting showed 4 students choosing Option B and one choosing Option A. When asked "Who would choose still chose B if this was your actual job tomorrow, with your manager evaluating you on engagement metrics?", no students switched from B to A. When asking the student who chose A, why they chose like that, they answered pragmatically. As a designer they want the best quality possible with modern technologies to provide the best possible experience for the user. When asked "Who is responsible if you choose A and the planet suffers?", they responded: They don't think they are at fault for what energy is consumed and its rather the manufacturers responsibility. The limited sample size prevents generalizable conclusions, but the exercise
illustrates a key tension: even designers with clear environmental values must
navigate professional contexts where sustainability competes with engagement
metrics. The one student who chose Option A articulated this tension: designers
are hired to provide optimal user experience with available technology, and
manufacturers bear responsibility for infrastructure efficiency. While this reflects legitimate professional reality, it underestimates designer
responsibility for two reasons. First, design choices create demand that
manufacturers must meet: responsibility is shared across the production chain.
Second, "optimal experience" doesnt need to mean maximum resource consumption. Many
users could prefer HD with faster loading if given informed choice. The structural pressure identified is real: designers who consistently prioritize
environmental impact over business metrics risk career consequences. This tension
between individual ethics and organizational incentives demonstrates why systemic
solutions like regulatory standards and mandatory disclosure are necessary complements
to individual responsibility.

\subsection{Possible Objections}
Having established that UX designers bear shared responsibility within structural constraints, three objections to this thesis warrant examination.
\begin{itemize}
  \item \textbf{Scale:} Data centers represent only 1.5\% of global emissions \cite{IEA2025EnergyAI}. Why prioritize UX ethics over larger sectors? This misunderstands the argument's scope. The claim is not that UX causes climate catastrophe, but that professionals with disproportionate influence over a rapidly growing source bear responsibility commensurate with their power. Moreover, dismissing every sector below 5\% of emissions would mean addressing none-climate change results from aggregating many "small" sources.

\item \textbf{Efficiency:} Won't technological improvements solve this without changing UX patterns? Google documented a 44-fold emissions reduction per AI query over 12 months, demonstrating rapid progress \cite{google}. However, total AI energy consumption continues growing \cite{IEA2025EnergyAI}. This pattern is characteristic of Jevons Paradox: efficiency improvements reduce
per-unit costs, which enables expanded consumption, resulting in higher total
resource use. The 19th-century economist William Stanley Jevons observed this
with coal. As engines became more efficient, coal consumption increased rather
than decreased because efficiency made coal-powered applications economically
viable in new contexts \cite{ALCOTT20059}. It has to be noted that recent scholarship questions whether Jevons Paradox applies straightforwardly to
complex AI systems \cite{jev}, but the pattern remains: without consumption
constraints, efficiency enables expansion rather than reduction.

  \item \textbf{Business:} Won't mandatory disclosure and default changes harm businesses? The objection treats profit as inviolable while treating environmental harm as negotiable. This is the inverse of the justice framework established in Section III. When regulation applies equally to competitors, it shifts industry norms rather than disadvantaging individual actors. This Objection also showcases that regulations must happen on a global scale.
\end{itemize}

\par Designers face genuine dilemmas between personal values and professional
incentives. Individual ethics alone cannot overcome structural barriers that
systematically reward engagement over sustainability. Section 5 therefore proposes
systemic interventions that enable ethical design to become professionally viable.
